\section{The marked block integral on the original activity addresses}
\label{sec:marked-block-integral}
We now obtain the exact tilted identity used in
Section~\ref{sec:leaf-expansion}. The insertion is transported through
the same integral and the same connected-component decomposition as
the background. There is no division by a forced activity, which may
vanish, and the distinguished activity address remains fixed.

\subsection{The fine inserted density}
Write the fine density as
\[
 \rho(U)=e^{E(U)}\sum_{A\ {\rm admissible}}\prod_{i\in A}K_i(U).
\]
Choose a bounded real measurable gauge-invariant function $h$, with
$|h|\le w$, supported on the eligibility condition of one fixed
activity address. Its dependence on the configuration is confined to
that address's full inspected support. There are two possibilities.
For a regular address $b$, its background regular activity is
identically zero and $h$ vanishes outside the healthy set for its
support. For a forced address $b$, $h$ vanishes whenever that forced
address fails its original eligibility test. Define
\begin{equation}\label{eq:marked-fine-density}
 f_h(U)=e^{E(U)}\sum_{A\ {\rm admissible}} B_h(A,U),\qquad
 B_h(A,U)=
 \begin{cases}
 h(U)\prod_{i\in A}K_i(U),&b\text{ regular},\\
 \mathbf1_{\{b\in A\}}h(U)\prod_{i\in A\setminus\{b\}}K_i(U),
                                      &b\text{ forced}.
 \end{cases}
\end{equation}
The admissibility tests, including exact flag coverage, are those
of the background. In the forced case the address is replaced even
when its original value $K_b(U)$ is zero.

Here and below use the same funded sector data and all-species
integrated product bounds as in
Section~\ref{sec:gaussian-reassembly}. On the fixed finite lattice,
choose summable uniform prices $p_i\ge|K_i(U)|$. In the regular case
the absolute family sum in \eqref{eq:marked-fine-density} is bounded by
$w\prod_i(1+p_i)$. In the forced case replace $p_b$ by $w$ and use
the same product bound. Thus the series converges uniformly and
absolutely. The fine exponent has a bounded measurable envelope, so
$f_h$ is bounded and measurable. Gauge invariance follows termwise
from the invariance of $h$, the activities and the flag tests.
In particular the fine Haar integral and every fixed transition
slice integral exist. This argument does not require a sign for $h$.

\subsection{One marked factor in the Gaussian expansion}
Denote the elementary alphabet of the unmarked expansion by
$\mathcal I$ and the letter assigned to $b$ by $\star$. Its support
is the same enlarged inspected support as the original address.
For one sector, one tag assignment and one old Haar configuration,
let $v_i(z)$ be the elementary factors of the Gaussian expansion.
The marked factor is
\[
 H(z)=h(U(z))\,\zeta_\star(z),\qquad 0\le\zeta_\star\le1,
\]
where $\zeta_\star$ is the larger local fluctuation cutoff on the
marked support. For a finite set of letters $P$ put
\begin{equation}\label{eq:marked-letter-product}
 \Psi_h(P,z)=\mathbf1_{\{\star\in P\}}H(z)
                          \prod_{i\in P\setminus\{\star\}}v_i(z).
\end{equation}
The cutoff makes the marked factor part of the same local integral
construction. On the nonzero support of the genuine common chart
cutoff it equals one: the stronger common small-fluctuation bound
places each marked coordinate in its plateau. If the common cutoff
is zero, both sides of the chart identity are zero. Hence inserting
$\zeta_\star$ changes none of the represented fine integral.

The regular and forced addresses give two slightly different
expansion identities. At a blank regular address the old
exponential-minus-one factor $v_\star$ is zero. Every old term
containing that letter therefore vanishes. Insert $\star$ into each
remaining finite letter set and replace its zero factor by $H$.
Erasing $\star$ is the inverse map, so the resulting sum is $H$
times the original sum. In the forced case, a family not containing
the address contributes zero. For a family containing it, erase
just that forced factor and put $H$ in its place. The optional
logarithmic and warning expansions are unchanged. These two
identities give precisely the two cases of
\eqref{eq:marked-fine-density}; there is no extra multiplicity.

To justify integration of these expansions, let $r_i\ge0$ be the
summable elementary integrated prices already established for the
unmarked product construction. Its bound applies to every finite
set of letters, in particular to $P\setminus\{\star\}$. Since
$|H|\le w$, it gives
\begin{equation}\label{eq:marked-integrated-price}
 \int|\Psi_h(P,z)|\,d\nu(z)
 \le \mathbf1_{\{\star\in P\}}w
                            \prod_{i\in P\setminus\{\star\}}r_i,
\end{equation}
where $\nu$ includes the old probability Haar variables and the
Gaussian variables. The right side is summable over $P$: erasure
maps the sets containing $\star$ injectively to all finite sets,
whose product-price sum is at most $\exp(\sum_i r_i)$.
The separate uniform envelope prices give measurability and an
integrable envelope exactly as in the unmarked construction.
Thus Fubini and absolute summation justify every exchange here.

For clarity, the chart and Gaussian normalization in this argument
are unchanged. First split the free Haar coordinates into the
active chart coordinates and their complement. On a nonzero sector
the chart-support lemma puts every active bond in the original
exponential chart. The product change of variables contributes
$c_H^{N_{\rm active}}\prod J(u_i)$; the affine rescaling contributes
$g^{qN_{\rm active}}$; the unused Gaussian coordinates integrate to
one with their normalized density. The quadratic completion and
logarithmic cancellation from
Section~\ref{sec:gaussian-reassembly} then apply with the extra
factor $h(U)$ untouched. Its local cutoff has just been restored
to one on this support. Summing the sector partition and averaging
the original tags consequently recovers $\mathcal T f_h$, with the
same scalar overhead $\mathscr O$ as for $\mathcal T\rho$.

\subsection{The distinguished connected component}
Split a finite elementary letter set into its connected components,
using their full assigned supports. If it contains $\star$, exactly
one component $C$ contains that letter. Formula
\eqref{eq:marked-letter-product} factors into the marked product on
$C$ and the ordinary products on the other components. The Gaussian
cross covariances between their read sets vanish; the Haar and tag
read sets are disjoint. Thus the product integral factors. This is
the same finite-range independence argument used in the unmarked
case, now with one bounded local factor substituted.

Let $a_C^F(V)$ be an ordinary integrated component value and
$b_C^F(V)$ the marked value, set to zero unless $\star\in C$.
The component and flattening maps are inverse on the permitted
families. The sector and eligibility masks are identical on both
sides: if a component fails its local test, its corresponding
ordinary or marked value is zero; on a nonzero family the union of
its inventories is exactly the set of flags. Absolute convergence
from \eqref{eq:marked-integrated-price} therefore proves
\begin{equation}\label{eq:marked-forest-integral}
 \mathcal T f_h(V)=\mathscr O
 \sum_{\mathcal A\ {\rm admissible}}
    \sum_{C\in\mathcal A}b_C^F(V)
                       \prod_{D\in\mathcal A\setminus\{C\}}a_D^F(V).
\end{equation}

The marked value has the same locality in the forest as the ordinary
one. Indeed $h$ reads only its inspected support, and its added
cutoff reads only the corresponding buffered coordinates. Forests
agreeing there give equal values after pulling their free Haar
variables back to the common full Haar space. Their principal
Gaussian covariances agree, so integration gives equality too.
We may therefore average independent forest choices at coarse
sites as in Section~\ref{sec:forest-average-log}. For a compatible
family the marked component and all the ordinary components use
disjoint sets of site choices, so the average of their product is
the product of their averages. Set
\[
             a_C=\mathbb E_F a_C^F,\qquad b_C=\mathbb E_F b_C^F.
\]
Equation~\eqref{eq:marked-forest-integral} holds with these averaged
values. The same absolute majorants apply to this finite average.

Two distinct compatible components cannot both have nonzero marked
value, since both would contain $\star$. Hence, for every real
$\theta$ and every compatible family,
\[
 \prod_{C\in\mathcal A}(a_C+\theta b_C)
 =\prod_{C\in\mathcal A}a_C
  +\theta\sum_{C\in\mathcal A}b_C
                           \prod_{D\ne C}a_D.
\]
The terms of degree at least two vanish individually. Together
with the unmarked identity this proves the exact integral formula
\begin{equation}\label{eq:marked-affine-integral}
 \mathcal T\rho+\theta\mathcal T f_h
   =\mathscr O\sum_{\mathcal A\ {\rm admissible}}
                               \prod_{C\in\mathcal A}(a_C+\theta b_C).
\end{equation}
The right side is affine despite its product notation. The masks,
support assignments, inventories and distinguished fine address
have never changed.

\subsection{Passing the mark through the two regroupings}
For $\theta$ in the interval allowed by the purchased marked prices,
apply the regular cluster expansion and forced exclusion regrouping
to the values $a_C+\theta b_C$. Denote their regular exponential
by $a(\theta)$ and their forced gas by $G_\theta$. The exact
identities of Sections~\ref{sec:forest-average-log}
and~\ref{sec:forced-exclusion} give
\[
 \mathcal T\rho+\theta\mathcal T f_h
                         =\mathscr O a(\theta)G_\theta.
\]
The absolute product and cluster majorants apply uniformly on this
interval. Indeed if $|a_C|\le A_C$, $|b_C|\le D B_C$ and
$|\theta|D\le1$, the affine factor has the joint value/derivative
bound $A_C+2B_C$ of Section~\ref{sec:leaf-expansion}. Products,
cluster coefficients and exponential multiplicities retain that
joint bound by the product rule proved there. This uses the
marked-price purchases; the unmarked prices alone do not supply
the derivative bound.

Finally let $R_0,S_0$ be the two fixed local logarithmic corrections
absorbed into the forced gas in Section~\ref{sec:soft-log-fusion}.
That regrouping uses the same fixed background and gives
$b(\theta)=G_\theta e^{R_0+S_0}$. Its background normalization is
$e^{E_{\rm coarse}+R_0+S_0}=\mathscr O a(0)$.
Substitution now proves
\[
 \mathcal T\rho+\theta\mathcal T f_h
 =e^{E_{\rm coarse}}\frac{a(\theta)}{a(0)}b(\theta),
\]
which is \eqref{eq:leaf-tilted-identity} for the actual insertion
\eqref{eq:marked-fine-density}. Differentiation is justified by
the joint summable bounds, not by a formal manipulation of the
series. This establishes the exact marked integral underlying that
argument. The remaining quantitative task is to purchase its marked
local derivative prices and principal-term estimate on the selected
history; the density identity alone does not provide those estimates.
